% ============================================================================
\section{Experimental Setup}
\label{sec:experiments}
% ============================================================================

We study a single Qwen3-1.7B student and four specialist teachers for
mathematics, code, instruction following, and science. Each task has one fixed
teacher. Appendix Table~\ref{tab:pipeline} lists the complete model, data,
optimization, and system metadata; entries marked in red remain to be filled
from the final run configuration.

\subsection{Data, loss, and fixed budget}

Each micro-batch contains four prompts from one task, each with one sampled
response, and is processed in one backward pass. A step contains 16
micro-batches and therefore 64 responses. Reverse KL is averaged over valid
tokens within each response and then over the four responses; $w_i/m_i$ is
applied only after this internal mean. Every attempted response contributes to
the response count and resource total, including invalid or truncated outputs.

The objective weights are the normalized inverse initial task losses. Every run
starts from the same student and uses these weights unchanged. The first step
uses Uniform and initializes the online statistics; later steps follow the
assigned count rule. With 500 steps, the total budget is $32{,}000$ attempted
responses. The learning-rate schedule uses optimizer step as its clock.

Each task's fixed stream contains $16{,}000$ prompts. Even a task assigned the
maximum count on every step consumes exactly
$500\times8\times4=16{,}000$ prompts, while Uniform consumes $8{,}000$.
Consequently, no prompt is reused. At a given step, methods can differ in the
number of prompts consumed from each task-specific shuffled stream.

\subsection{Methods and focused comparisons}

The primary comparison uses conventional AdamW for Uniform, GPAS, and
Cost-GPAS, isolating the count rule from the optimizer. Two focused diagnostics
address the method's central design choices: RawNoise replaces the
AdamW-scaled noise by raw-coordinate noise, and a matched Uniform/GPAS pair uses
the optional taskwise (TW) second moment. The count bounds, signal definitions,
and second-moment choices are fixed before the reported training loss is
examined. Uniform assigns $(4,4,4,4)$ micro-batches. GPAS uses
$w_i\sqrt{e_i}$, and Cost-GPAS minimizes
Equation~\ref{eq:batch_cost_objective}. RawNoise substitutes unscaled gradient
noise for $e_i$. Uniform--TW and GPAS--TW use the corresponding count rules
with the optional second-moment observation in
Equation~\ref{eq:taskwise_square_observation}. Appendix
Table~\ref{tab:exp_methods} gives a compact summary.

All configurations consume the same shuffled prompt stream for a task in the
same order: the $n$th requested prompt is identical across configurations.
This pairing removes prompt-order variation from method comparisons. The
allocation determines how many prompts each task has consumed at a given step.

\subsection{Research questions and evaluation}

The experiments answer five questions:
\begin{enumerate}
    \item How different are the tasks' optimizer-scaled noise levels, how do
    they change during training, and when does raw noise rank them differently?
    \item Does GPAS reduce held-out gradient variance and weighted teacher loss
    at a matched number of optimizer steps?
    \item Does Cost-GPAS reduce the GPU hours required to reach a common loss?
    \item Does the optional taskwise second moment materially change the result?
    \item Does every downstream task remain at least as strong as its Uniform
    counterpart under the fixed objective?
\end{enumerate}

At early, middle, and late checkpoints from Uniform, each task supplies a
fixed set of independently generated micro-batches. We estimate the allocation
on one half and measure its weighted-gradient variance on the other, then
exchange the halves. This paired test evaluates the allocation separately from
the online variance ratio $H$.

Capability evaluation includes MATH-500 greedy pass@1, a fixed LiveCodeBench
pass@1 slice, IFBench strict accuracy, and GPQA-Diamond average@4.
\missing{Confirm benchmark versions and decoding settings; add AIME24/25
average@16 or a second science benchmark if the listed set has insufficient
range above the initial student.} Alongside the raw scores, we average the
per-task normalized gain
$(s-s_{\rm init})/(s_{\rm teacher}-s_{\rm init})$, using the initial student and
the corresponding specialist teacher as endpoints. We also report average
response length and truncation rate by task.

\subsection{System and accounting}

At most one teacher is loaded on the GPU at a time. Each step loads all four
teachers in a fixed cyclic order, and every method uses this order because the
target hardware cannot hold all teacher parameters simultaneously. GPU hours
equal full wall time times the number of reserved GPUs, including model-transfer
and waiting time.
Each trace records teacher readiness, rollout, scoring, backward, optimizer
time, tokens, attempted responses, and peak memory. The serial timing model
extracts $\tau_i$ from marginal per-micro-batch time and $C$ from the remaining
fixed work. We report $C/(\sum_i m_i\tau_i)$ next to end-to-end GPU hours.
Appendix~\ref{app:system_details} gives the timing definitions and resume test.

% ============================================================================
\section{Results}
\label{sec:results}
% ============================================================================

\subsection{Noise and allocation dynamics}

\begin{figure}[H]
    \centering
    \missingfigure[0.18\textheight]{Insert Figure 2 from the four-teacher logs:
    (a) scaled and raw noise by task over training; (b) GPAS and Cost-GPAS
    counts; (c) variance ratio $H$ with the bound $2$; (d) task times and fixed
    overhead ratio.}
    \caption{Optimizer-scaled noise and the resulting allocations.
    \missing{Replace the placeholder with logged trajectories and state the
    main quantitative observation in the caption.}}
    \label{fig:allocation_dynamics}
\end{figure}

\missing{Report the cross-task range and temporal drift of $\widehat e_i$, the
fraction of steps where raw and scaled noise disagree in rank, how often each
task reaches a count bound, and the median and quantiles of $H$. If $H$ stays
near one, state that Uniform is sufficient in this regime.}

\subsection{Training efficiency}

\begin{figure}[H]
    \centering
    \missingfigure[0.17\textheight]{Insert Figure 3: (a) weighted teacher loss
    versus optimizer step for Uniform and GPAS; (b) weighted teacher loss versus
    GPU hours for Uniform, GPAS, and Cost-GPAS. Put per-task curves in the
    appendix.}
    \caption{Training progress at matched update and compute budgets.
    \missing{Replace the placeholder and add uncertainty intervals.}}
    \label{fig:training_efficiency}
\end{figure}

\missing{State the matched-step loss difference for GPAS, the GPU-hour saving
for Cost-GPAS at the predeclared Uniform-reachable threshold, and the associated
response/token counts. Base each comparison on its paired uncertainty interval.}

\subsection{Capability and resource summary}

\begin{table}[H]
\centering
\scriptsize
\setlength{\tabcolsep}{2.5pt}
\caption{Four-teacher capability and efficiency.
\missing{Populate from the fixed evaluation outputs and timing logs.}}
\label{tab:main_results}
\begin{tabular}{lccccccc}
\toprule
& Math & Code & IF & Science & Norm. avg. & Final $F$ & GPU h to target \\
\midrule
Initial student & & & & & & & \\
Teacher upper bounds & & & & & & & \\
Uniform & & & & & & & \\
GPAS & & & & & & & \\
Cost-GPAS & & & & & & & \\
\bottomrule
\end{tabular}
\end{table}

\missing{Report the initial student and specialist endpoints, all task scores,
normalized average, final weighted teacher loss, and GPU hours to threshold.
State every task-level difference from Uniform alongside the average.}

\subsection{Two targeted diagnostics}

\begin{table}[H]
\centering
\small
\caption{Allocation signal and optional optimizer-state diagnostics.
\missing{Fill after RawNoise and TW runs.}}
\label{tab:ablations}
\begin{tabular}{llccc}
\toprule
Count rule & AdamW state & Median $H$ & Final $F$ & Norm. avg. \\
\midrule
Uniform & Conv. & & & \\
RawNoise & Conv. & & & \\
GPAS & Conv. & & & \\
Uniform & TW & & & \\
GPAS & TW & & & \\
\bottomrule
\end{tabular}
\end{table}

\missing{Use the RawNoise comparison to isolate optimizer coordinates and the
matched Conv./TW pairs to decide whether the optional correction matters. Keep
other hyperparameters fixed; no further allocation ablations are needed.}

\subsection{Controlled calculation}

The controlled three-task check isolates the allocation arithmetic before the
large-model study. GPAS attains $0.679\times$ Uniform's optimizer-scaled
variance, close to the calculated $0.681\times$, whereas RawNoise reaches
$2.316\times$. Zero-overhead Cost-GPAS attains a time--variance product of
$0.857\times$ Uniform, close to the calculated $0.854\times$.
Appendix~\ref{app:controlled} contains the setup, figure, and second-moment
checks.
